% !TeX program = xelatex
% Самодостатній файл. Компілятор Overleaf: XeLaTeX.
% Розклад: титул 2 хв + 4 блоки по 17 хв = 70 хв.
% Усі числові приклади є навчальними моделями, а не реальними вимірюваннями.
\documentclass[aspectratio=169,11pt]{beamer}
\usepackage{fontspec}
\usepackage{polyglossia}
\setdefaultlanguage{ukrainian}
\setotherlanguage{english}
\IfFontExistsTF{DejaVu Sans}{\setmainfont{DejaVu Sans}\setsansfont{DejaVu Sans}\newfontfamily\cyrillicfont{DejaVu Sans}\newfontfamily\cyrillicfontsf{DejaVu Sans}}{\setmainfont{CMU Serif}\setsansfont{CMU Sans Serif}\newfontfamily\cyrillicfont{CMU Serif}\newfontfamily\cyrillicfontsf{CMU Sans Serif}}
\usepackage{amsmath,amssymb,mathtools,booktabs}
\usefonttheme{professionalfonts}
\definecolor{Ink}{HTML}{172F43}
\definecolor{Accent}{HTML}{006E79}
\definecolor{Muted}{HTML}{536471}
\setbeamercolor{normal text}{fg=Ink,bg=white}
\setbeamercolor{frametitle}{fg=Ink,bg=white}
\setbeamercolor{structure}{fg=Accent}
\setbeamercolor{footline}{fg=Muted,bg=white}
\setbeamerfont{frametitle}{size=\large,series=\bfseries}
\setbeamertemplate{navigation symbols}{}
\setbeamertemplate{headline}{}
\setbeamersize{text margin left=9mm,text margin right=9mm}
\setbeamertemplate{frametitle}{\vspace{2mm}\insertframetitle\par\vspace{1mm}}
\newcommand{\lecturelabel}{}
\newcommand{\blocklabel}{}
\newcommand{\slottime}{}
\setbeamertemplate{footline}{\hspace{9mm}\makebox[.64\paperwidth][l]{\scriptsize\lecturelabel\quad\blocklabel}\hfill{\scriptsize\slottime\quad\insertframenumber}\hspace{9mm}\vspace{3mm}}
\setlength{\parskip}{1.8mm}
\setlength{\parindent}{0pt}
\newcommand{\p}[1]{\par\vspace{1mm}{\color{Accent}\bfseries #1}\quad}
\newcommand{\source}[2]{\par\vspace{1mm}{\scriptsize\color{Muted}\href{#2}{#1}}}
\newcommand{\question}[1]{\par\vspace{2mm}{\color{Accent}#1}}
\hypersetup{unicode=true,colorlinks=true,urlcolor=Accent,linkcolor=Ink,pdftitle={Інформатика та інформаційні технології — 16 лекцій},pdfauthor={Навчальні матеріали за робочою програмою 2022 року}}
% За потреби нотатки можна показати заміною hide notes на show notes.
\setbeameroption{hide notes}
\begin{document}

% ==================== ЛЕКЦІЯ 2 ====================
\renewcommand{\lecturelabel}{Лекція 2}
\renewcommand{\blocklabel}{}
\section{2. Похибки і довіра до результату}
\renewcommand{\slottime}{2 хв}
\begin{frame}[t]{}
\vspace{7mm}
{\color{Accent}\large Лекція 2}\par\vspace{3mm}
{\LARGE\bfseries Похибки і довіра до результату}\par\vspace{4mm}
Скільки цифр у відповіді справді мають зміст?\par\vfill
{\small Інформатика та інформаційні технології\par 70 хвилин. Чотири навчальні блоки}

\note{Озвучити прикладне запитання лекції. Зібрати одну-дві короткі відповіді аудиторії й пов’язати їх із попередньою темою. Титульний слайд разом зі вступом займає 2 хвилини.}
\end{frame}


% Блок 2.1. 17 хв: 2 + 4 + 4 + 4 + 3.
\renewcommand{\blocklabel}{Блок 1}
\subsection{Джерела похибки}
\renewcommand{\slottime}{2 хв}
\begin{frame}[t]{Джерела похибки}
Неточність з’являється у вимірюваннях, спрощеннях моделі та обчисленнях. Їх треба розрізняти, щоб покращувати саме слабке місце.
\question{Чому багато знаків після коми не роблять оцінку ремонту достовірною?}
\vfill
\p{Екзаменаційні питання} 6.
\note{Пояснити ідею звичайними словами. Чому багато знаків після коми не роблять оцінку ремонту достовірною? Дати аудиторії коротко спрогнозувати відповідь до введення формул.}
\end{frame}

\renewcommand{\slottime}{4 хв}
\begin{frame}[t]{Застосування і цікаве спостереження}
\p{Де це потрібно} Рулетка дає похибку довжини. Модель прямокутної кімнати ігнорує ніші. Комп’ютер округлює числа під час обчислень.
\p{Цікаве спостереження} Число $0{,}1$ має нескінченний двійковий запис. Тому звичайний двійковий float зберігає лише його наближення.\source{Python: двійкова арифметика}{https://docs.python.org/3/tutorial/floatingpoint.html}
\note{Розгорнути перший приклад на слайді, попросити назвати вхідні дані та практичний результат. Пов’язати цікаве спостереження з обмеженнями методу. Джерело: Python: двійкова арифметика.}
\end{frame}

\renewcommand{\slottime}{4 хв}
\begin{frame}[t]{Формальне означення і формула}
Сумарну помилку можна розкласти через проміжні точні значення моделі та методу.
\[
|y_{\rm real}-\widehat y|\le e_{\rm model}+e_{\rm data}+e_{\rm method}+e_{\rm round}.
\]
\p{Умови й позначення} Це межа через нерівність трикутника. Реальна похибка може бути меншою через взаємну компенсацію складових.
\note{Послідовно пояснити кожне позначення та зв’язати його з прикладом. Проговорити умови застосування. Не вимагати запам’ятовування формули до розуміння її дії.}
\end{frame}

\renewcommand{\slottime}{4 хв}
\begin{frame}[t]{Розбір навчального прикладу}
Навчальний бюджет похибок площі: модель до 0,4 м$^2$, вимірювання до 0,2 м$^2$, обчислення до 0,01 м$^2$.
\[
e_{\rm total}\le0{,}4+0{,}2+0{,}01=0{,}61\ \text{м}^2.
\]
\question{Яку складову варто зменшити першою, якщо потрібна межа 0,3 м$^2$?}
\note{Виконати обчислення на дошці або усно з аудиторією. Спочатку запропонувати передбачити знак або порядок результату. Яку складову варто зменшити першою, якщо потрібна межа 0,3 м$^2$? Обговорити відповідь і перевірку одиниць, якщо вони є.}
\end{frame}

\renewcommand{\slottime}{3 хв}
\begin{frame}[t]{Ідея завдання на Python}
\p{Ідея Python} Бюджет похибок кошторису.\p{Дані} Три задані невід’ємні межі похибок.\p{Результат} Їх сума та найбільша складова.\p{Перевірка} Зменшення найменшої складової не обов’язково суттєво покращує загальний результат.
\note{Обговорити вхід, вихід і послідовність дій. Запропонувати один звичайний і один граничний приклад. На лекції не писати повну реалізацію. Пов’язати завершення блоку з наступною темою.}
\end{frame}


% Блок 2.2. 17 хв: 2 + 4 + 4 + 4 + 3.
\renewcommand{\blocklabel}{Блок 2}
\subsection{Абсолютна і відносна похибки}
\renewcommand{\slottime}{2 хв}
\begin{frame}[t]{Абсолютна і відносна похибки}
Абсолютна похибка вимірює відхилення в одиницях величини. Відносна показує, наскільки велике це відхилення порівняно з самою величиною.
\question{Однакова помилка 1 см є суттєвою для деталі й майже непомітною для довжини кімнати. Як це порівняти?}
\vfill
\p{Екзаменаційні питання} 7, 8.
\note{Пояснити ідею звичайними словами. Однакова помилка 1 см є суттєвою для деталі й майже непомітною для довжини кімнати. Як це порівняти? Дати аудиторії коротко спрогнозувати відповідь до введення формул.}
\end{frame}

\renewcommand{\slottime}{4 хв}
\begin{frame}[t]{Застосування і цікаве спостереження}
\p{Де це потрібно} Виготовлення вставки потребує допуску в міліметрах. Порівняння точності зважування різних вантажів зручне у відсотках.
\p{Цікаве спостереження} Поблизу нуля відносна похибка стає незручною: навіть мале абсолютне відхилення може дати великий відсоток.
\note{Розгорнути перший приклад на слайді, попросити назвати вхідні дані та практичний результат. Пов’язати цікаве спостереження з обмеженнями методу. Спостереження виводиться з формул і навчальних даних цього блоку.}
\end{frame}

\renewcommand{\slottime}{4 хв}
\begin{frame}[t]{Формальне означення і формула}
Для точного $x$ і наближеного $\widehat x$ абсолютна похибка дорівнює $\Delta$, її гарантована межа дорівнює $\Delta_{\max}$.
\[
\Delta=|x-\widehat x|,\quad \delta=\frac{\Delta}{|x|};\qquad \delta\le\frac{\Delta_{\max}}{|\widehat x|-\Delta_{\max}}.
\]
\p{Умови й позначення} Відносну похибку визначено при $x\ne0$. Остання межа справедлива при $|\widehat x|>\Delta_{\max}$ і не потребує знання $x$.
\note{Послідовно пояснити кожне позначення та зв’язати його з прикладом. Проговорити умови застосування. Не вимагати запам’ятовування формули до розуміння її дії.}
\end{frame}

\renewcommand{\slottime}{4 хв}
\begin{frame}[t]{Розбір навчального прикладу}
Еталонна довжина 100 мм, виміряне значення 99 мм. Для іншої деталі еталон 10 мм, результат 9 мм.
\[
\Delta_1=\Delta_2=1\ \text{мм},\qquad \delta_1=1\%,\quad\delta_2=10\%.
\]
\question{Який прилад обрати, якщо важливий допуск 0,5 мм, а не відсоток?}
\note{Виконати обчислення на дошці або усно з аудиторією. Спочатку запропонувати передбачити знак або порядок результату. Який прилад обрати, якщо важливий допуск 0,5 мм, а не відсоток? Обговорити відповідь і перевірку одиниць, якщо вони є.}
\end{frame}

\renewcommand{\slottime}{3 хв}
\begin{frame}[t]{Ідея завдання на Python}
\p{Ідея Python} Перевірка вимірювання деталі.\p{Дані} Еталон, вимір і допустима абсолютна похибка.\p{Результат} Абсолютна похибка, відносна похибка та відповідність допуску.\p{Перевірка} Для нульового еталона відносну похибку не обчислювати.
\note{Обговорити вхід, вихід і послідовність дій. Запропонувати один звичайний і один граничний приклад. На лекції не писати повну реалізацію. Пов’язати завершення блоку з наступною темою.}
\end{frame}


% Блок 2.3. 17 хв: 2 + 4 + 4 + 4 + 3.
\renewcommand{\blocklabel}{Блок 3}
\subsection{Значущі цифри та округлення}
\renewcommand{\slottime}{2 хв}
\begin{frame}[t]{Значущі цифри та округлення}
Значущі цифри починаються з першої ненульової цифри. Точність запису має відповідати точності даних, а правило округлення повинно бути визначеним.
\question{Чи однаково інформативні записи 2,5 м і 2,500 м?}
\vfill
\p{Екзаменаційні питання} 9.
\note{Пояснити ідею звичайними словами. Чи однаково інформативні записи 2,5 м і 2,500 м? Дати аудиторії коротко спрогнозувати відповідь до введення формул.}
\end{frame}

\renewcommand{\slottime}{4 хв}
\begin{frame}[t]{Застосування і цікаве спостереження}
\p{Де це потрібно} Вимірювання маси зразка й запис результату в протокол. Округлення загальної суми покупки до копійок.
\p{Цікаве спостереження} Python для точних половин використовує округлення до найближчого парного: наприклад, round(2.5) дає 2. Це відрізняється від шкільного правила «половину вгору».\source{Python: правило round}{https://docs.python.org/3/library/functions.html\#round}
\note{Розгорнути перший приклад на слайді, попросити назвати вхідні дані та практичний результат. Пов’язати цікаве спостереження з обмеженнями методу. Джерело: Python: правило round.}
\end{frame}

\renewcommand{\slottime}{4 хв}
\begin{frame}[t]{Формальне означення і формула}
Округлення до найближчого кратного кроку $u$ має похибку не більше $u/2$. Вірність знаків тут розуміємо у вузькому сенсі.
\[
|x-\operatorname{round}_u(x)|\le\frac u2,\qquad \Delta_{\max}\le\tfrac12\,10^{m-n+1}.
\]
\p{Умови й позначення} У другій формулі $10^m\le|\widehat x|<10^{m+1}$, а $n$ є числом гарантовано вірних значущих цифр за цим критерієм.
\note{Послідовно пояснити кожне позначення та зв’язати його з прикладом. Проговорити умови застосування. Не вимагати запам’ятовування формули до розуміння її дії.}
\end{frame}

\renewcommand{\slottime}{4 хв}
\begin{frame}[t]{Розбір навчального прикладу}
Для 12,347 при округленні до 0,01 маємо 12,35. Межа похибки округлення 0,005, фактична похибка 0,003.
\[
|12{,}347-12{,}35|=0{,}003\le0{,}005.
\]
\question{Чому округлення кожного доданка може дати іншу суму, ніж округлення загального результату?}
\note{Виконати обчислення на дошці або усно з аудиторією. Спочатку запропонувати передбачити знак або порядок результату. Чому округлення кожного доданка може дати іншу суму, ніж округлення загального результату? Обговорити відповідь і перевірку одиниць, якщо вони є.}
\end{frame}

\renewcommand{\slottime}{3 хв}
\begin{frame}[t]{Ідея завдання на Python}
\p{Ідея Python} Порівняння двох способів округлення цін.\p{Дані} Кілька цін з трьома десятковими знаками.\p{Результат} Сума округлених цін та округлена сума.\p{Перевірка} Окремо розглянути 2.5 і 3.5 та назвати правило Python.
\note{Обговорити вхід, вихід і послідовність дій. Запропонувати один звичайний і один граничний приклад. На лекції не писати повну реалізацію. Пов’язати завершення блоку з наступною темою.}
\end{frame}


% Блок 2.4. 17 хв: 2 + 4 + 4 + 4 + 3.
\renewcommand{\blocklabel}{Блок 4}
\subsection{Поширення похибок}
\renewcommand{\slottime}{2 хв}
\begin{frame}[t]{Поширення похибок}
Якщо результат залежить від кількох вимірювань, неточність кожного з них впливає на відповідь. Чутливі аргументи варто вимірювати точніше.
\question{Для оцінки площі столу важливіша точність довжини чи ширини?}
\vfill
\p{Екзаменаційні питання} 10.
\note{Пояснити ідею звичайними словами. Для оцінки площі столу важливіша точність довжини чи ширини? Дати аудиторії коротко спрогнозувати відповідь до введення формул.}
\end{frame}

\renewcommand{\slottime}{4 хв}
\begin{frame}[t]{Застосування і цікаве спостереження}
\p{Де це потрібно} Площа матеріалу $S=ab$. Густина зразка $\rho=m/V$: малий об’єм може посилити вплив похибки його вимірювання.
\p{Цікаве спостереження} При подвоєнні радіуса площа круга зростає вчетверо. Для малих похибок відносна похибка площі приблизно удвічі більша за похибку радіуса.
\note{Розгорнути перший приклад на слайді, попросити назвати вхідні дані та практичний результат. Пов’язати цікаве спостереження з обмеженнями методу. Спостереження виводиться з формул і навчальних даних цього блоку.}
\end{frame}

\renewcommand{\slottime}{4 хв}
\begin{frame}[t]{Формальне означення і формула}
Лінеаризоване поширення похибок для диференційовної функції $f(x_1,\ldots,x_n)$ задається її частинними похідними.
\[
\Delta f\ \lesssim\ \sum_{i=1}^{n}\left|\frac{\partial f}{\partial x_i}\right|\Delta_i.
\]
\p{Умови й позначення} Це оцінка першого порядку для малих збурень. Строга межа використовує максимуми модулів похідних на всій області можливих даних.
\note{Послідовно пояснити кожне позначення та зв’язати його з прикладом. Проговорити умови застосування. Не вимагати запам’ятовування формули до розуміння її дії.}
\end{frame}

\renewcommand{\slottime}{4 хв}
\begin{frame}[t]{Розбір навчального прикладу}
Нехай $a=2\pm0{,}01$ м, $b=3\pm0{,}01$ м. Для добутку можна також врахувати член другого порядку.
\[
\Delta S\le3\cdot0{,}01+2\cdot0{,}01+0{,}01^2=0{,}0501\ \text{м}^2.
\]
\question{Яке вимірювання доцільно уточнити, якщо точність можна покращити лише для однієї сторони?}
\note{Виконати обчислення на дошці або усно з аудиторією. Спочатку запропонувати передбачити знак або порядок результату. Яке вимірювання доцільно уточнити, якщо точність можна покращити лише для однієї сторони? Обговорити відповідь і перевірку одиниць, якщо вони є.}
\end{frame}

\renewcommand{\slottime}{3 хв}
\begin{frame}[t]{Ідея завдання на Python}
\p{Ідея Python} Калькулятор меж площі.\p{Дані} Дві додатні довжини та їхні менші за них похибки.\p{Результат} Номінальна площа, мінімальна й максимальна можливі площі.\p{Перевірка} Порівняти точний інтервал з лінеаризованою оцінкою.
\note{Обговорити вхід, вихід і послідовність дій. Запропонувати один звичайний і один граничний приклад. На лекції не писати повну реалізацію. Пов’язати завершення блоку з наступною темою.}
\end{frame}


\end{document}
